% Chapter 3: Gödel's Incompleteness Theorems
% Part I: The Nature of Computation

\chapter{Gödel's Incompleteness Theorems}
\label{ch:godel}

%======================================================================
% SECTION 1: WHAT IS A FORMAL SYSTEM?
%======================================================================
\section{What Is a Formal System?}
\label{sec:godel:formalsystem}

By the early 20th century, mathematics faced a foundational crisis. Paradoxes in set theory (Russell's paradox, 1901) showed that intuitive reasoning about infinite collections leads to contradictions. Hilbert's program (1920s) proposed a solution: formalize all of mathematics in a precise axiomatic system, then prove using \emph{finitary} methods that this system is:

\begin{enumerate}
\item \textbf{Consistent:} No contradiction can be derived ($P \land \neg P$ is unprovable).
\item \textbf{Complete:} Every true statement is provable (for any $\phi$, either $\phi$ or $\neg\phi$ is provable).
\item \textbf{Decidable:} There is an algorithm to determine whether any given statement is provable (the Entscheidungsproblem).
\end{enumerate}

Can mathematics prove every truth? Can a formal system prove its own consistency?

Hilbert believed the answer was yes. He famously declared in 1930: \emph{``Wir müssen wissen, wir werden wissen''} (``We must know, we shall know''). Gödel proved the answer is \emph{no} --- twice. To understand why, we first need to understand precisely what a formal system is, and why a whole generation of mathematicians believed that one could be built to contain all of mathematics.

\subsection{The Anatomy of a Formal System}
\label{sec:godel:anatomy}

What, exactly, was the object Hilbert wanted to construct? A \emph{formal system} is a mathematical object made of pure syntax --- a precise, rule-governed language with no meaning attached:

\begin{definition}[Formal System]
A formal system consists of:
\begin{enumerate}
\item An \textbf{alphabet}: a finite set of symbols (e.g., $0, S, +, \times, =, \neg, \to, \forall, (, )$ and variables).
\item \textbf{Formation rules}: rules for which finite strings of symbols are well-formed \emph{formulas}.
\item \textbf{Axioms}: a finite list (or finite schema) of formulas taken as starting points.
\item \textbf{Rules of inference}: mechanical rules for deriving new formulas from old ones (e.g., modus ponens: from $\phi$ and $\phi \to \psi$, derive $\psi$).
\end{enumerate}
A \emph{proof} (derivation) is a finite sequence of formulas, each of which is either an axiom or is obtained from earlier formulas by a rule of inference. A \emph{theorem} is the last formula of some proof.
\end{definition}

The defining feature is that a formal system is \emph{mechanical}: each proof step is a finite, checkable manipulation of symbols, requiring no insight. This is why Hilbert could compare operating with formulas to playing chess --- the rules, not the meaning, do all the work. Peano Arithmetic (PA) is the canonical example: its axioms are Robinson's arithmetic $\mathsf{Q}$ together with the induction schema, and its only rule of inference is modus ponens. Gödel's theorems apply to any formal system that is \emph{consistent} and \emph{strong enough} to express elementary arithmetic. Nothing about such a system, in itself, refers to truth or meaning; those live only in the \emph{interpretation} we give to the symbols.

\subsection{Hilbert's Program: The Three Pillars}
\label{sec:godel:hilbert}

The second of Hilbert's famous 23 problems (1900) asked for a proof of the consistency of the arithmetic axioms. Hilbert believed that such a proof would forever secure mathematics against paradox. The program had two phases: \emph{axiomatization} (express all mathematics in formal language) and \emph{consistency proving} (use finitary metamathematics to show no contradiction can be derived). When he retired in 1930, his address to the Society of German Scientists and Physicians reiterated his unwavering faith: no mathematical truth would remain unknown. Within a year, a 25-year-old Austrian logician had shattered this vision.

The idea that a formal system could be both complete and consistent was not just hoped for --- it was \emph{assumed}. The alternative was unthinkable:

\begin{enumerate}
\item \textbf{The Intuition of Completeness:} Mathematics discovers eternal truths. A proper formalization should capture \emph{all} of them. Incompleteness would mean mathematics has ``blind spots'' --- true statements that are formally unprovable. This seemed like a defect in the formalism, not in mathematics itself.

\item \textbf{The Intuition of Consistency:} If mathematics is the study of necessary truths, it cannot contain contradictions. But proving consistency \emph{within} the system seemed circular --- you would be using the system to vouch for itself. Hilbert's hope was to use \emph{weaker} (finitary) methods to prove the consistency of stronger systems.
\end{enumerate}

The prevailing view was that any foundational crisis could be resolved by more careful axiomatization. After all, Zermelo-Fraenkel set theory (1908--1922) had eliminated the known paradoxes. Hilbert's program was the logical next step: prove once and for all that the new foundations were sound.

\subsection{The Pre-Gödel Landscape}
\label{sec:godel:landscape}

Before 1931, several approaches tried to secure foundations:

\begin{itemize}
\item \textbf{Logicism (Frege, Russell):} Reduce mathematics to logic. Crushed by Russell's paradox (1901). \emph{Principia Mathematica} tried to patch with type theory, introducing a hierarchy of types to ban self-membership.
\item \textbf{Formalism (Hilbert):} Treat mathematics as symbol manipulation. Prove consistency finitistically. This was the dominant view in the 1920s. Hilbert believed all mathematical problems were solvable in principle.
\item \textbf{Intuitionism (Brouwer):} Reject actual infinity and the law of excluded middle. Constructive mathematics only. Rejected by most mainstream mathematicians as too restrictive, though it anticipated important ideas in computability.
\item \textbf{Zermelo-Fraenkel Set Theory (1908, 1922):} Axiomatize set theory. Still needed a consistency proof. The axiom of choice remained controversial.
\end{itemize}

Gödel's 1929 Completeness Theorem (every valid first-order formula is provable) raised hopes --- but it applied to \emph{first-order logic}, not to arithmetic with induction. The completeness theorem gave a perfect match between semantic validity and syntactic provability for first-order logic. If arithmetic were similarly complete, Hilbert's program would succeed.

\subsection{The Road to 1931}
\label{sec:godel:road}

The key developments that converged in Gödel's breakthrough:

\begin{itemize}
\item \textbf{1901:} Russell's Paradox destroys Frege's logicism.
\item \textbf{1908:} Russell's Type Theory; Zermelo's set theory axioms.
\item \textbf{1910--13:} \emph{Principia Mathematica} (Whitehead \& Russell) --- 2000 pages to derive $1+1=2$.
\item \textbf{1920s:} Hilbert develops metamathematics (proof theory) as a discipline.
\item \textbf{1929:} Gödel's \emph{Completeness Theorem} for first-order logic: every valid formula is provable. This gave hope --- maybe \emph{arithmetic} is complete too?
\item \textbf{1930:} Gödel's PhD (Completeness). At the Königsberg conference, he quietly announces: ``I have a result that shows Hilbert's program is impossible.''
\item \textbf{1931:} ``Über formal unentscheidbare Sätze der Principia Mathematica und verwandter Systeme I'' --- the Incompleteness Theorems.
\end{itemize}

\subsection{Gödel's Life and Context}
\label{sec:godel:life}

Kurt Friedrich Gödel was born April 28, 1906 in Brünn, Austria-Hungary (now Brno, Czech Republic). A shy, precocious child, he excelled in mathematics and languages. He entered the University of Vienna in 1924, originally intending to study physics, but was drawn to mathematics and logic.

Vienna in the 1920s was a vibrant intellectual center. The \textbf{Vienna Circle} (Wiener Kreis), led by Moritz Schlick, advocated logical positivism: all knowledge is either empirical or analytic (true by definition). Gödel attended their meetings but never fully subscribed to their philosophy --- his Platonism (belief in the reality of mathematical objects) set him apart from the Circle's empiricism. Yet the Circle's emphasis on logical analysis and the structure of scientific theories shaped his approach.

The \textbf{Königsberg Conference} of September 1930 was a turning point. At a meeting on epistemology and the exact sciences, Gödel presented his completeness theorem. During a roundtable discussion on Hilbert's program --- attended by Carnap, Heyting, and von Neumann --- Gödel offhandedly mentioned that he had found a way to construct formally undecidable propositions. Von Neumann immediately grasped the significance and later corresponded with Gödel. \cite{godel1931}. By 1931, the full paper was published in \emph{Monatshefte für Mathematik und Physik}.

The response was dramatic. Hilbert initially reacted with anger, then grudging acceptance. Bernays, Hilbert's collaborator, worked with Gödel to refine the derivability conditions. The logical community split: some saw incompleteness as the end of foundational dreams; others (including Gödel himself) saw it as opening new avenues.

Gödel later moved to the Institute for Advanced Study in Princeton (1933 onward), where he became close friends with Albert Einstein. The two walked home together daily from the Institute. Einstein remarked that he went to his office ``just for the privilege of walking home with Gödel.'' Their conversations ranged from physics to philosophy to politics. Gödel's Platonism influenced Einstein's later views on the nature of mathematical reality.

Gödel's later work included the consistency of the Continuum Hypothesis with ZFC (the constructible universe $L$, 1938), contributions to relativity theory (Gödel metric, 1949 --- a rotating universe solution allowing closed timelike curves), and deep philosophical writings on the nature of mathematics. In 1948, he became a US citizen; at his citizenship hearing, he attempted to explain to the judge how the US Constitution could be amended to establish a dictatorship (Einstein and Princeton colleagues had to steer him back on track). He became increasingly paranoid in his later years, fearing poisoning and believing that others were conspiring against him. He refused to eat, convinced that his food was poisoned, and eventually starved himself to death on January 14, 1978, with a weight of only 65 pounds.

Gödel's philosophical legacy is as significant as his mathematical one. He wrote extensively on the philosophy of mathematics, defending Platonism against both formalism and intuitionism. His 1951 Gibbs Lecture, ``Some Basic Theorems on the Foundations of Mathematics,'' argued that the incompleteness theorems decisively refute the view that mathematics is a purely formal game. He also contributed to Leibniz scholarship, believing that Leibniz had discovered a universal characteristic (``characteristica universalis'') that could resolve all disputes through calculation --- a vision that anticipates modern formal verification and AI reasoning.

%======================================================================
% SECTION 2: SELF-REFERENCE
%======================================================================
\section{Self-Reference}
\label{sec:godel:selfref}

The key to the whole affair was a phenomenon everyone else regarded as a pathology: \emph{self-reference}. A sentence that talks about itself seemed to be precisely the kind of thing that breaks logic. Gödel's masterstroke was to treat self-reference not as a bug to be eliminated but as an instrument to be engineered --- by smuggling it into arithmetic, where no one expected to find it. This section develops the idea; the encoding machinery that makes it fully rigorous is developed in Section~\ref{sec:godel:numbering}.

\subsection{The Liar Paradox}
\label{sec:godel:liar}

``This statement is false.'' If true, it's false. If false, it's true. Classical logic bans such statements. The paradox is ancient --- Epimenides, a Cretan, reportedly declared that all Cretans are liars --- but its sharp modern form is the self-referential liar sentence $L$: ``$L$ is false.'' Whatever truth value one assigns to $L$, one contradicts oneself. The difficulty is not the sentence's subject matter; it is the sentence's ability to talk about \emph{its own truth value}. Any language that can express self-reference seems to contain a logical time bomb.

\subsection{Self-Reference as a Bug}
\label{sec:godel:selfref_bug}

Before Gödel, self-reference in formal systems was considered a bug to be eliminated, not a tool. Russell's type theory (1908) stratified propositions to prevent self-reference: a sentence may only quantify over types strictly below its own, so no sentence can refer to itself or to its own truth. Tarski (1933) later showed that a language cannot define its own truth predicate without paradox --- truth in a language can only be defined in a richer meta-language, never in the language itself. The lesson drawn from the tradition was clear: self-reference breaks logic; keep it out.

But the intuition that self-reference \emph{must} be avoided was about to be overturned. Gödel realized that self-reference could be \emph{encoded} in the innocuous language of arithmetic, where no one expected it to live.

\subsection{Gödel's Leap: Harnessing Self-Reference}
\label{sec:godel:leap}

Statements \emph{about} a formal system (provability, consistency) can be translated into statements \emph{within} the system (arithmetic) by encoding symbols, formulas, and proofs as numbers. The statement ``Formula $\phi$ is provable'' becomes an arithmetic statement about the Gödel number of $\phi$. Then, using diagonalization, construct a sentence $G$ that says ``I am not provable.'' If $G$ is provable, the system is inconsistent. If $G$ is not provable, $G$ is true but unprovable --- the system is incomplete.

This leap had three audacious components:

\begin{enumerate}
\item \textbf{Arithmetization of Syntax:} Every symbol $\to$ number. Every formula $\to$ number (via prime factorization). Every proof $\to$ number. The relation ``$x$ is the Gödel number of a proof of the formula with Gödel number $y$'' becomes a primitive recursive predicate $\mathsf{Proof}(x, y)$.

\item \textbf{The Fixed Point (Diagonal) Lemma:} For any formula $\phi(x)$ with one free variable, there exists a sentence $G$ such that the system proves $G \leftrightarrow \phi(\ulcorner G \urcorner)$. Applied to $\phi(x) = \neg \mathsf{Prov}(x)$, we get $G \leftrightarrow \neg \mathsf{Prov}(\ulcorner G \urcorner)$.

\item \textbf{Truth $\neq$ Provability:} The sentence $G$ is \emph{true} (in the standard model $\N$) because it asserts its own unprovability, and indeed it is unprovable (assuming consistency). But the system cannot prove $G$. Truth has escaped provability.
\end{enumerate}

Before Gödel, ``true'' and ``provable'' were conflated for formal systems. Hilbert's program \emph{required} them to coincide (completeness). Gödel showed they diverge: there are arithmetic truths that no finite proof can reach. The map is not the territory. More subtly, Gödel showed that \emph{every} consistent, sufficiently strong system suffers this divergence --- it is not a defect of a particular system but a law of mathematical reality.

\subsection{The Structure of the Proof}
\label{sec:godel:structure}

The proof is a masterpiece of encoding and self-reference:

\begin{enumerate}
\item \textbf{Encoding:} Map every syntactic object (symbol, formula, proof) to a natural number (its Gödel number).
\item \textbf{Primitive Recursiveness:} Show that syntactic relations (``$x$ is a formula,'' ``$x$ proves $y$'') are primitive recursive, hence representable in PA.
\item \textbf{Diagonalization:} Construct a sentence that refers to its own Gödel number using the fixed-point lemma.
\item \textbf{Truth in $\N$:} The sentence says ``I am not provable.'' If the system is consistent, it is true but unprovable.
\item \textbf{Consistency $\to$ Incompleteness:} If the system could prove its own consistency, it would prove $G$ --- contradiction.
\end{enumerate}

\subsection{The Diagonal Lemma}
\label{sec:godel:diagonal}

\begin{lemma}[Diagonal / Fixed Point Lemma]
For any formula $\phi(v)$ with one free variable, there exists a sentence $G$ such that:
\[
\mathsf{PA} \vdash G \leftrightarrow \phi(\ulcorner G \urcorner)
\]
\end{lemma}

\begin{proof}[Proof Sketch]
Define $\mathsf{diag}(n) = \ulcorner \phi(\ulcorner \phi \urcorner) \urcorner$ where $n = \ulcorner \phi \urcorner$. This function is recursive. Let $\delta(v) = \phi(\mathsf{diag}(v))$. Let $d = \ulcorner \delta \urcorner$. Then $G = \delta(d)$ satisfies $G \leftrightarrow \phi(\ulcorner G \urcorner)$.
\end{proof}

\subsubsection*{Detailed Construction of the Fixed Point}

Step 1: Define the substitution function $\mathsf{subst}(n, m)$. If $n = \ulcorner \phi(v) \urcorner$ (a formula with one free variable), then $\mathsf{subst}(n, m) = \ulcorner \phi(\overline{m}) \urcorner$ (the result of substituting the numeral for $m$ for the free variable in $\phi$). This function is primitive recursive --- substitution is a purely syntactic operation.

Step 2: Define the diagonal function $\mathsf{diag}(n) = \mathsf{subst}(n, n)$. If $n = \ulcorner \phi(v) \urcorner$, then $\mathsf{diag}(n) = \ulcorner \phi(\ulcorner \phi(v) \urcorner) \urcorner = \ulcorner \phi(\overline{n}) \urcorner$.

Step 3: Given $\phi(v)$, define $\delta(v) := \phi(\mathsf{diag}(v))$. Let $d = \ulcorner \delta \urcorner$. Then:
\[
G := \delta(d) = \phi(\mathsf{diag}(d))
= \phi(\ulcorner \delta(d) \urcorner)
= \phi(\ulcorner G \urcorner)
\]

So PA proves $G \leftrightarrow \phi(\ulcorner G \urcorner)$ because $\mathsf{diag}(d)$ computes $\ulcorner G \urcorner$ by construction.

The crucial insight: $G$ is literally the same syntactic object as $\phi(\ulcorner G \urcorner)$ after computation of the diagonal function. Self-reference is achieved not by semantic magic but by syntactic construction.

%======================================================================
% SECTION 3: GÖDEL NUMBERING
%======================================================================
\section{Gödel Numbering}
\label{sec:godel:numbering}

Self-reference is achieved not by magic but by arithmetic. The trick is \emph{Gödel numbering}: turn the symbols, formulas, and proofs of a formal system into natural numbers. Then the system --- which talks about numbers --- can talk about its own syntax. This is the arithmetization of metamathematics: the deepest philosophical questions about mathematics become questions about the arithmetic of natural numbers.

\subsection{Encoding Mathematics}
\label{sec:godel:encoding}

Let $\mathcal{L}$ be the language of Peano Arithmetic (PA): symbols $\{0, S, +, \times, =, \neg, \to, \forall, (, ), x_1, x_2, \dots\}$.

\begin{definition}[Gödel Numbering]
Assign each symbol a unique natural number. A sequence of symbols $s_1 s_2 \dots s_k$ is encoded as:
\[
\ulcorner s_1 s_2 \dots s_k \urcorner = 2^{n_1} \cdot 3^{n_2} \cdot 5^{n_3} \cdots p_k^{n_k}
\]
where $p_i$ is the $i$-th prime ($p_1 = 2, p_2 = 3, p_3 = 5, \dots$) and $n_i$ is the code of $s_i$. By unique prime factorization, decoding is effective.
\end{definition}

\subsubsection*{A Concrete Encoding Table}

A typical assignment of Gödel numbers to the basic symbols:

\[
\begin{array}{c|c}
\text{Symbol} & \text{Gödel number} \\ \hline
0 & 1 \\
S & 2 \\
= & 3 \\
+ & 4 \\
\times & 5 \\
\neg & 6 \\
\to & 7 \\
\forall & 8 \\
( & 9 \\
) & 10 \\
x_k & 11 + k \text{ (for variable index } k \text{)}
\end{array}
\]

\subsubsection*{Worked Example: Encoding a Formula}

Consider the simple formula $\forall x_1 (x_1 = x_1)$. The symbol sequence is:

\[
\forall,\; x_1,\; (,\; x_1,\; =,\; x_1,\; )
\]

The Gödel numbers are: $8,\; 12,\; 9,\; 12,\; 3,\; 12,\; 10$.

The Gödel number of the entire formula is:

\[
2^{8} \cdot 3^{12} \cdot 5^{9} \cdot 7^{12} \cdot 11^{3} \cdot 13^{12} \cdot 17^{10}
\]

This is a huge number, but it is a well-defined natural number. Decoding it recovers the original formula uniquely.

For the sentence $0 = 0$:

\[
\ulcorner 0 = 0 \urcorner = 2^{1} \cdot 3^{3} \cdot 5^{1} = 2 \cdot 27 \cdot 5 = 270
\]

Even simple formulas produce surprisingly large codes because of the exponential growth from prime powers.

\subsubsection*{Encoding Sequences of Formulas (Proofs)}

A proof is a finite sequence of formulas $\phi_1, \phi_2, \dots, \phi_n$ where each $\phi_i$ is either an axiom or follows from earlier formulas by a rule of inference. The Gödel number of the proof is:

\[
\ulcorner \phi_1 \dots \phi_n \urcorner = 2^{\ulcorner \phi_1 \urcorner} \cdot 3^{\ulcorner \phi_2 \urcorner} \cdot 5^{\ulcorner \phi_3 \urcorner} \cdots p_n^{\ulcorner \phi_n \urcorner}
\]

This is a double exponential encoding, but it is still a fully effective coding.

\subsection{Primitive Recursive Functions}
\label{sec:godel:primrec}

To represent syntactic relations inside PA, we need a class of functions simple enough that their definitions can be expressed in the language of arithmetic, yet expressive enough to capture all mechanical syntactic manipulation.

\begin{definition}[Primitive Recursive Functions]
The class of primitive recursive functions is the smallest class containing:
\begin{enumerate}
\item The \textbf{zero function}: $Z(x) = 0$ for all $x$.
\item The \textbf{successor function}: $S(x) = x+1$.
\item The \textbf{projection functions}: $P_i^n(x_1, \dots, x_n) = x_i$ for $1 \le i \le n$.
\end{enumerate}
and closed under:
\begin{enumerate}
\item[4.] \textbf{Composition}: If $g$ and $h_1, \dots, h_m$ are primitive recursive, then $f(x_1, \dots, x_n) = g(h_1(\vec{x}), \dots, h_m(\vec{x}))$ is primitive recursive.
\item[5.] \textbf{Primitive Recursion}: If $g$ (arity $n$) and $h$ (arity $n+2$) are primitive recursive, then the function $f$ defined by:
\[
\begin{aligned}
f(x_1, \dots, x_n, 0) &= g(x_1, \dots, x_n) \\
f(x_1, \dots, x_n, y+1) &= h(x_1, \dots, x_n, y, f(x_1, \dots, x_n, y))
\end{aligned}
\]
is primitive recursive.
\end{enumerate}
\end{definition}

\subsubsection*{Examples of Primitive Recursive Functions}

\begin{itemize}
\item \textbf{Addition:} $\mathsf{add}(x, 0) = x$, $\mathsf{add}(x, y+1) = S(\mathsf{add}(x, y))$.
\item \textbf{Multiplication:} $\mathsf{mult}(x, 0) = 0$, $\mathsf{mult}(x, y+1) = \mathsf{add}(\mathsf{mult}(x, y), x)$.
\item \textbf{Exponentiation:} $x^0 = 1$, $x^{y+1} = \mathsf{mult}(x^y, x)$.
\item \textbf{Predecessor:} $\mathsf{pred}(0) = 0$, $\mathsf{pred}(y+1) = y$.
\item \textbf{Subtraction (truncated):} $x \dot- 0 = x$, $x \dot- (y+1) = \mathsf{pred}(x \dot- y)$.
\item \textbf{Factorial:} $0! = 1$, $(y+1)! = (y+1) \cdot y!$.
\end{itemize}

\subsubsection*{Primitive Recursive Predicates}

A predicate $P(\vec{x})$ is primitive recursive if its characteristic function $\chi_P(\vec{x})$ (which is 1 if $P$ holds, 0 otherwise) is primitive recursive.

Key primitive recursive predicates include:
\begin{itemize}
\item $x = y$ (equality)
\item $x \le y$ (order)
\item $x \mid y$ (divisibility)
\item $\mathsf{Prime}(x)$ (primality test)
\item $\mathsf{Even}(x)$ (parity)
\end{itemize}

The crucial point: all syntactic properties (``$x$ is a term,'' ``$x$ is a formula,'' ``$x$ is an axiom,'' ``$x$ is a proof of $y$'') are primitive recursive predicates. This is because syntactic manipulation is purely finite combinatorial manipulation of symbols.

\subsubsection*{Building the Syntactic Predicates Step by Step}

We can construct the key primitive recursive predicates systematically:

\begin{enumerate}
\item $\mathsf{Term}(x)$: $x$ is the Gödel number of a term. Terms are built from $0$, variables, and closure under $S$, $+$, $\times$. Since the language is finite, the predicate is defined by a primitive recursive search over term construction trees --- the depth of the term bounds the search.

\item $\mathsf{Formula}(x)$: $x$ is the Gödel number of a formula. Formulas are built from atomic formulas ($t_1 = t_2$) using $\neg$, $\to$, and $\forall$. Again, definable by primitive recursion on formula complexity.

\item $\mathsf{Axiom}(x)$: $x$ is the Gödel number of an axiom of PA. This includes the axioms of Robinson arithmetic $\mathsf{Q}$ (successor, addition, multiplication axioms), the induction schema, and logical axioms. Checking that a formula matches a finite list of schemas is primitive recursive.

\item $\mathsf{Proof}(x, y)$: $x$ is the Gödel number of a finite sequence of formulas $x = \ulcorner \phi_1, \dots, \phi_n \urcorner$ such that each $\phi_i$ is either an axiom or follows from earlier $\phi_j, \phi_k$ by modus ponens or generalization. This is primitive recursive because we can decode the sequence, iterate through the formulas, and check each condition using bounded search.
\end{enumerate}

\subsubsection*{The $\beta$-Function: Encoding Sequences as Numbers}

A technical challenge: we need to represent arbitrary finite sequences of natural numbers as single natural numbers. Gödel's solution was the $\beta$-function:

\begin{definition}[Gödel's $\beta$-function]
\[
\beta(a, b, i) = \text{the remainder when } a \text{ is divided by } 1 + (i+1) \cdot b
\]
\end{definition}

The remarkable property: for any finite sequence $c_0, c_1, \dots, c_k$, there exist numbers $a, b$ such that $\beta(a, b, i) = c_i$ for all $i \le k$. This relies on the Chinese Remainder Theorem and lets us quantify over sequences without using set theory.

\begin{definition}[Representability]
A function $f: \N^k \to \N$ is \textbf{representable} in PA if there exists a formula $\phi_f(x_1, \dots, x_k, y)$ such that:
\begin{itemize}
\item If $f(n_1, \dots, n_k) = m$, then $\mathsf{PA} \vdash \phi_f(\overline{n_1}, \dots, \overline{n_k}, \overline{m})$.
\item $\mathsf{PA} \vdash \exists! y\, \phi_f(\overline{n_1}, \dots, \overline{n_k}, y)$.
\end{itemize}
\end{definition}

\begin{theorem}[Representability of Primitive Recursive Functions]
Every primitive recursive function is representable in PA. Moreover, the representing formula can be chosen to be $\Sigma_1$ (a block of existential quantifiers followed by a quantifier-free formula).
\end{theorem}

This theorem is the bridge between syntax-as-numbers and the internal expressive power of PA. Because all syntactic predicates are primitive recursive, they are representable, which means PA can ``talk about'' its own syntax.

\subsection{The Provability Predicate}
\label{sec:godel:provpred}

Define the primitive recursive predicate:
\begin{itemize}
\item $\mathsf{Axiom}(x)$: $x$ is the Gödel number of an axiom of PA.
\item $\mathsf{InferenceRule}(x, y, z)$: $z$ follows from $x$ and $y$ by modus ponens (or generalization).
\item $\mathsf{ProofLength}(k, x, y)$: the sequence with Gödel number $x$ is a proof of length $k$ of the formula with Gödel number $y$.
\item $\mathsf{Proof}(x, y)$: $x$ is a proof of $y$.
\end{itemize}

\begin{definition}[Provability Predicate]
\[
\mathsf{Prov}(y) \equiv \exists x\, \mathsf{Proof}(x, y)
\]
\end{definition}

This is a $\Sigma_1$ formula --- it says ``there exists a natural number that encodes a proof of the formula with Gödel number $y$.'' When $y = \ulcorner \phi \urcorner$, $\mathsf{Prov}(\ulcorner \phi \urcorner)$ asserts ``$\phi$ is provable in PA.''

%======================================================================
% SECTION 4: THE INCOMPLETENESS THEOREMS
%======================================================================
\section{The Incompleteness Theorems}
\label{sec:godel:incompleteness}

All the machinery is now in place. With the encoding of Section~\ref{sec:godel:numbering} and the diagonal lemma of Section~\ref{sec:godel:selfref}, the two theorems follow almost directly.

\subsection{The First Incompleteness Theorem}
\label{sec:godel:first}

\begin{theorem}[Gödel's First Incompleteness Theorem (1931)]
Any consistent, recursively axiomatized extension of Robinson Arithmetic $\mathsf{Q}$ (or PA) is incomplete: there exists a sentence $G$ such that neither $G$ nor $\neg G$ is provable.
\end{theorem}

\begin{proof}
Apply the Diagonal Lemma to $\phi(v) = \neg \mathsf{Prov}(v)$ to obtain $G$ such that $\mathsf{PA} \vdash G \leftrightarrow \neg \mathsf{Prov}(\ulcorner G \urcorner)$.

Assume $\mathsf{PA}$ is consistent.

\begin{itemize}
\item Suppose $\mathsf{PA} \vdash G$. Then there is a proof of $G$, so $\mathsf{Prov}(\ulcorner G \urcorner)$ is true in $\N$. Since $\mathsf{Prov}$ is representable, $\mathsf{PA} \vdash \mathsf{Prov}(\ulcorner G \urcorner)$. But from $G \leftrightarrow \neg \mathsf{Prov}(\ulcorner G \urcorner)$, we have $\mathsf{PA} \vdash \neg \mathsf{Prov}(\ulcorner G \urcorner)$. Thus $\mathsf{PA}$ proves both $\mathsf{Prov}(\ulcorner G \urcorner)$ and $\neg \mathsf{Prov}(\ulcorner G \urcorner)$, making it inconsistent. Contradiction.

\item Suppose $\mathsf{PA} \vdash \neg G$. Then $\mathsf{PA} \vdash \mathsf{Prov}(\ulcorner G \urcorner)$ (by $G \leftrightarrow \neg \mathsf{Prov}(\ulcorner G \urcorner)$). So $\exists x\, \mathsf{Proof}(x, \ulcorner G \urcorner)$ holds in $\N$. If the system is $\omega$-consistent (as Gödel originally assumed), this implies there is a specific $n$ with $\mathsf{Proof}(n, \ulcorner G \urcorner)$, and thus $\mathsf{PA} \vdash G$. This gives $\mathsf{PA} \vdash G \land \neg G$, inconsistency.
\end{itemize}

If the system is merely consistent (not $\omega$-consistent), the second case requires Rosser's strengthening. But under $\omega$-consistency, neither $G$ nor $\neg G$ is provable.

$G$ is \emph{true} in the standard model $\N$: it asserts its own unprovability, and indeed it is unprovable (assuming consistency). Truth and provability diverge.
\end{proof}

\begin{definition}[$G$ is a \emph{Gödel Sentence}]
$G$ says ``I am not provable.'' It is true but unprovable (assuming consistency).
\end{definition}

The first theorem can be stated more concisely: \emph{For any consistent, recursively axiomatizable theory $T$ that interprets $\mathsf{Q}$, $T$ is incomplete.}

\subsubsection*{Incompleteness and the Halting Problem}

There is a deep connection between Gödel's first incompleteness theorem and the undecidability of the halting problem (Chapter~2). In fact, the incompleteness theorem can be derived from the undecidability of the halting problem:

Suppose PA were complete. Then for any Turing machine $M$, we could decide whether $M$ halts: encode the statement ``$M$ halts'' as an arithmetic sentence $\phi_M$, and then (assuming completeness) either $\phi_M$ or $\neg \phi_M$ is provable. By searching proofs systematically, we would eventually find the answer. This would make the halting problem decidable. Since the halting problem is undecidable, PA must be incomplete.

Conversely, the undecidability of the halting problem follows from the first incompleteness theorem: if the halting problem were decidable, we could decide provability (by checking if a proof search halts), contradicting incompleteness.

The two results are equivalent in the sense that each implies the other. They are the same phenomenon viewed through different lenses: a limit on formal systems (logical lens) and a limit on computation (computability lens).

\subsection{The Second Incompleteness Theorem}
\label{sec:godel:second}

\begin{definition}[Consistency Statement]
$\mathsf{Con}(\mathsf{PA}) \equiv \neg \mathsf{Prov}(\ulcorner 0=1 \urcorner)$ --- ``There is no proof of a contradiction.''
\end{definition}

\begin{theorem}[Gödel's Second Incompleteness Theorem (1931)]
If $\mathsf{PA}$ is consistent, then $\mathsf{PA} \nvdash \mathsf{Con}(\mathsf{PA})$.
\end{theorem}

\begin{proof}[Proof Sketch]
The proof of the First Theorem can be formalized \emph{inside} PA:
\[
\mathsf{PA} \vdash \mathsf{Con}(\mathsf{PA}) \to G
\]
(If PA is consistent, then $G$ is not provable, so $G$ holds.) But we know $\mathsf{PA} \nvdash G$ (if consistent). Therefore $\mathsf{PA} \nvdash \mathsf{Con}(\mathsf{PA})$.

The formalization requires that PA can prove the \emph{derivability conditions} (Hilbert-Bernays-Löb):
\begin{enumerate}
\item If $\mathsf{PA} \vdash \phi$, then $\mathsf{PA} \vdash \mathsf{Prov}(\ulcorner \phi \urcorner)$.
\item $\mathsf{PA} \vdash \mathsf{Prov}(\ulcorner \phi \to \psi \urcorner) \to (\mathsf{Prov}(\ulcorner \phi \urcorner) \to \mathsf{Prov}(\ulcorner \psi \urcorner))$.
\item $\mathsf{PA} \vdash \mathsf{Prov}(\ulcorner \phi \urcorner) \to \mathsf{Prov}(\ulcorner \mathsf{Prov}(\ulcorner \phi \urcorner) \urcorner)$.
\end{enumerate}
\end{proof}

\subsubsection*{The Derivability Conditions in Detail}

The three Hilbert-Bernays-Löb derivability conditions are the key to proving the second theorem:

\begin{enumerate}
\item \textbf{Internalization of Proof}: If the system proves $\phi$, then it also proves that it proves $\phi$. This requires that the proof predicate $\mathsf{Proof}$ is correctly defined.

\item \textbf{Closure under Modus Ponens}: The system proves that if $\phi \to \psi$ is provable and $\phi$ is provable, then $\psi$ is provable. This is the formalization of modus ponens inside the system.

\item \textbf{Internalization of Condition 1}: The system proves that if $\phi$ is provable, then it is provably provable. This is a reflection principle: the system can iterate its own provability predicate.
\end{enumerate}

Using these conditions, the formalization of the first theorem proceeds as follows:

From $G \leftrightarrow \neg \mathsf{Prov}(\ulcorner G \urcorner)$, we have:
\[
\mathsf{PA} \vdash \mathsf{Prov}(\ulcorner G \urcorner) \to \mathsf{Prov}(\ulcorner \neg \mathsf{Prov}(\ulcorner G \urcorner) \urcorner)
\]

By condition 3, $\mathsf{PA} \vdash \mathsf{Prov}(\ulcorner G \urcorner) \to \mathsf{Prov}(\ulcorner \mathsf{Prov}(\ulcorner G \urcorner) \urcorner)$.

Using condition 2, $\mathsf{PA} \vdash \mathsf{Prov}(\ulcorner G \urcorner) \to \mathsf{Prov}(\ulcorner \mathsf{Prov}(\ulcorner G \urcorner) \land \neg \mathsf{Prov}(\ulcorner G \urcorner) \urcorner)$.

Thus $\mathsf{PA} \vdash \mathsf{Prov}(\ulcorner G \urcorner) \to \mathsf{Prov}(\ulcorner 0=1 \urcorner)$ (since a contradiction proves anything).

Contrapositively: $\mathsf{PA} \vdash \neg \mathsf{Prov}(\ulcorner 0=1 \urcorner) \to \neg \mathsf{Prov}(\ulcorner G \urcorner)$.

But $\neg \mathsf{Prov}(\ulcorner G \urcorner)$ is equivalent to $G$ (by the fixed point). So:
\[
\mathsf{PA} \vdash \mathsf{Con}(\mathsf{PA}) \to G
\]

If $\mathsf{PA} \vdash \mathsf{Con}(\mathsf{PA})$, then $\mathsf{PA} \vdash G$, contradicting the first theorem. Hence $\mathsf{PA} \nvdash \mathsf{Con}(\mathsf{PA})$.

\subsubsection*{Löb's Theorem (1955)}

Löb's theorem generalizes the second incompleteness theorem and provides deep insight into the nature of provability:

\begin{theorem}[Löb's Theorem]
If $\mathsf{PA} \vdash \mathsf{Prov}(\ulcorner \phi \urcorner) \to \phi$, then $\mathsf{PA} \vdash \phi$.
\end{theorem}

\begin{proof}[Proof Sketch]
Apply the diagonal lemma to $\psi(v) := \mathsf{Prov}(v) \to \phi$ to get $L \leftrightarrow (\mathsf{Prov}(\ulcorner L \urcorner) \to \phi)$. Using the derivability conditions, one can show $\mathsf{PA} \vdash L \to \phi$ and $\mathsf{PA} \vdash \mathsf{Prov}(\ulcorner L \urcorner) \to \mathsf{Prov}(\ulcorner \phi \urcorner)$. From the assumption $\mathsf{Prov}(\ulcorner \phi \urcorner) \to \phi$, we get $\mathsf{PA} \vdash \mathsf{Prov}(\ulcorner L \urcorner) \to \phi$, so $\mathsf{PA} \vdash L$, hence $\mathsf{PA} \vdash \phi$.
\end{proof}

Löb's theorem shows that the system cannot prove its own soundness for any non-theorem. Taking $\phi = \bot$ (contradiction), we get: if $\mathsf{PA} \vdash \neg \mathsf{Prov}(\ulcorner \bot \urcorner)$, then $\mathsf{PA} \vdash \bot$. This is exactly the second incompleteness theorem.

Löb's theorem also implies that the provability predicate $\mathsf{Prov}$ satisfies the modal logic GL (Gödel-Löb), whose defining axiom is $\Box(\Box p \to p) \to \Box p$. This modal logic captures the structural behavior of provability in any sufficiently strong formal system.

Löb's theorem says: the only formulas $\phi$ for which the system can prove ``if $\phi$ is provable then $\phi$ is true'' are those that are already theorems. The system cannot prove its own soundness on any non-theorem. This is a precise formalization of the idea that the system cannot ``trust'' itself beyond what it has already proven.

\subsection{Why Truth Differs From Proof}
\label{sec:godel:truth_vs_proof}

The first incompleteness theorem's central phenomenon is the divergence of \emph{truth} from \emph{provability}. The sentence $G$ is true in the standard model $\N$ --- it asserts its own unprovability, and indeed it is unprovable --- yet the system cannot prove it. Why should this be?

The reason is that the two notions live at different logical levels. Provability is a $\Sigma_1$ property: $\mathsf{Prov}(y) \equiv \exists x\, \mathsf{Proof}(x, y)$, where the existence is witnessed by a concrete number --- the code of a proof. Truth, by contrast, is not definable in the language at all. Tarski's theorem makes this precise:

\begin{theorem}[Tarski's Undefinability of Truth]
There is no formula $\mathsf{True}(v)$ in the language of PA such that for every sentence $\phi$:
\[
\mathsf{PA} \vdash \mathsf{True}(\ulcorner \phi \urcorner) \leftrightarrow \phi
\]
\end{theorem}

\begin{proof}
Suppose such a $\mathsf{True}$ exists. Apply the diagonal lemma to $\neg \mathsf{True}(v)$ to obtain a sentence $L$ such that $\mathsf{PA} \vdash L \leftrightarrow \neg \mathsf{True}(\ulcorner L \urcorner)$. But by the property of $\mathsf{True}$, we have $\mathsf{PA} \vdash \mathsf{True}(\ulcorner L \urcorner) \leftrightarrow L$. Combining, $\mathsf{PA} \vdash L \leftrightarrow \neg L$, a contradiction. Therefore no such truth predicate exists.
\end{proof}

Tarski's theorem is the formal version of the liar paradox: ``This statement is false'' cannot be expressed within the language without contradiction. While Gödel showed that \emph{provability} is expressible (as the $\Sigma_1$ predicate $\mathsf{Prov}$), Tarski showed that \emph{truth} (in the standard model) is not. Note the contrast: $\mathsf{Prov}$ is a $\Sigma_1$ formula that captures provability, but it does not capture truth --- there are true but unprovable sentences. Truth is a $\Pi_1^1$ notion (truth in the standard model requires quantifying over all subsets of $\N$), far beyond the expressive power of a $\Sigma_1$ formula.

There is also a subtlety worth confronting directly: does Gödel's 1929 \emph{completeness} theorem contradict the incompleteness theorem? There is no contradiction: completeness applies to first-order \emph{logic} (the deductive apparatus), while incompleteness applies to specific \emph{theories} (like PA). The logical rules are complete: everything that is logically valid is provable. But the axioms of PA do not capture all arithmetic truths --- the set of models of PA includes non-standard models (with infinite numbers) where different sentences are true.

This distinction is crucial: completeness of logic does not imply completeness of theories. Gödel himself was surprised by this --- he had hoped that the completeness theorem would help prove the consistency of the Continuum Hypothesis, but his incompleteness results blocked that approach.

\subsection{Rosser's Improvement (1936)}
\label{sec:godel:rosser}

Gödel's theorem assumed \emph{$\omega$-consistency} (a stronger condition than consistency). Rosser removed this by constructing a sentence $R$ that says: ``For any proof of me, there is a shorter proof of my negation.''

\begin{definition}[$\omega$-consistency]
A theory $T$ is $\omega$-consistent if there is no formula $\phi(x)$ such that $T \vdash \exists x\, \phi(x)$ but $T \vdash \neg \phi(\overline{n})$ for every natural number $n$.
\end{definition}

$\omega$-consistency implies consistency but is strictly stronger. Rosser showed that mere consistency suffices for incompleteness.

\begin{theorem}[Rosser's Theorem]
If PA is consistent (not $\omega$-consistent), then PA is incomplete.
\end{theorem}

\begin{proof}[Proof]
Let $\mathsf{Proof}(x, y)$ mean ``$x$ is the Gödel number of a proof of the formula with Gödel number $y$.'' Define the Rosser provability predicate:

\[
\mathsf{RosserProv}(y) \equiv \exists x\, [\mathsf{Proof}(x, y) \land \forall z < x\; \neg \mathsf{Proof}(z, \neg y)]
\]

where $\neg y$ denotes $\ulcorner \neg \phi \urcorner$ when $y = \ulcorner \phi \urcorner$. This says ``there is a proof of $\phi$ such that no shorter proof proves $\neg \phi$.''

Apply the diagonal lemma to $\phi(v) = \neg \mathsf{RosserProv}(v)$ to obtain $R$ such that:
\[
\mathsf{PA} \vdash R \leftrightarrow \neg \mathsf{RosserProv}(\ulcorner R \urcorner)
\]

Now assume PA is consistent. We show neither $R$ nor $\neg R$ is provable.

\textbf{Case 1:} Suppose $\mathsf{PA} \vdash R$. Then there is a proof of $R$; let $n$ be the Gödel number of the shortest such proof. For all $k < n$, $k$ is not a proof of $\neg R$ (otherwise $\mathsf{PA}$ would prove $\neg R$, making it inconsistent). Thus $\mathsf{RosserProv}(\ulcorner R \urcorner)$ holds, so $\mathsf{PA} \vdash \mathsf{RosserProv}(\ulcorner R \urcorner)$. But $R \leftrightarrow \neg \mathsf{RosserProv}(\ulcorner R \urcorner)$ gives $\mathsf{PA} \vdash \neg \mathsf{RosserProv}(\ulcorner R \urcorner)$. Contradiction.

\textbf{Case 2:} Suppose $\mathsf{PA} \vdash \neg R$. Then there is a proof of $\neg R$; let $m$ be the Gödel number of the shortest such proof. For all $k < m$, $k$ is not a proof of $R$ (otherwise $\mathsf{PA}$ would prove both $R$ and $\neg R$). Moreover, $m$ itself is not a proof of $R$, so there is \emph{no} proof of $R$ that has no shorter proof of $\neg R$. Hence $\neg \mathsf{RosserProv}(\ulcorner R \urcorner)$ holds, so $\mathsf{PA} \vdash \neg \mathsf{RosserProv}(\ulcorner R \urcorner)$. But then $R \leftrightarrow \neg \mathsf{RosserProv}(\ulcorner R \urcorner)$ gives $\mathsf{PA} \vdash R$, contradicting consistency.

Therefore neither $R$ nor $\neg R$ is provable: PA is incomplete (without assuming $\omega$-consistency).
\end{proof}

\subsection{Gentzen's Consistency Proof (1936) and Proof-Theoretic Ordinals}
\label{sec:godel:gentzen}

Gentzen (1936) proved $\mathsf{Con}(\mathsf{PA})$ using transfinite induction up to $\varepsilon_0$, the first fixed point of $\alpha \mapsto \omega^\alpha$:
\[
\varepsilon_0 = \sup \{ \omega, \omega^\omega, \omega^{\omega^\omega}, \dots \}
\]

This does \emph{not} violate the second incompleteness theorem because transfinite induction up to $\varepsilon_0$ is \emph{not} formalizable in PA. The proof-theoretic ordinal of PA is $\varepsilon_0$ --- it measures the strength of PA's induction.

\begin{definition}[Proof-Theoretic Ordinal]
The proof-theoretic ordinal of a theory $T$ is the supremum of ordinals $\alpha$ such that $T$ proves transfinite induction up to $\alpha$ (for primitive recursive well-orderings).
\end{definition}

\begin{itemize}
\item Proof-theoretic ordinal of PA: $\varepsilon_0$
\item Proof-theoretic ordinal of $\mathsf{ATR}_0$ (arithmetical transfinite recursion): $\Gamma_0$ (Feferman-Schütte ordinal)
\item Proof-theoretic ordinal of ZFC: unknown (far beyond $\Gamma_0$)
\end{itemize}

Gentzen's proof uses \emph{cut elimination} for a sequent calculus formulation of PA. Each cut elimination step reduces the ordinal assigned to the proof, and the process terminates because there are no infinite descending sequences of ordinals below $\varepsilon_0$.

\subsubsection*{Ordinal Analysis in Detail}

The ordinal $\varepsilon_0$ can be represented in Cantor normal form:
\[
\alpha = \omega^{\beta_1} + \omega^{\beta_2} + \dots + \omega^{\beta_k}
\]
with $\beta_1 \ge \beta_2 \ge \dots \ge \beta_k$. The ordering is lexicographic on the exponents.

Gentzen assigns ordinals to proofs:
\begin{itemize}
\item Axioms get ordinal 0.
\item A cut with cut formula of complexity $\alpha$ increases the ordinal.
\item Cut elimination reduces the ordinal.
\end{itemize}

The proof-theoretic ordinal of a theory is the least ordinal that the theory \emph{cannot} prove is well-ordered. For PA, this is $\varepsilon_0$.

\begin{theorem}[Gentzen 1936]
$\mathsf{PA} + \mathsf{TI}(\varepsilon_0) \vdash \mathsf{Con}(\mathsf{PA})$, where $\mathsf{TI}(\varepsilon_0)$ is transfinite induction up to $\varepsilon_0$.
\end{theorem}

\subsubsection*{Beyond PA: Ordinal Analysis}

\begin{itemize}
\item \textbf{Robinson Arithmetic $\mathsf{Q}$}: $\omega$ (very weak)
\item \textbf{Second-order arithmetic $\mathsf{ACA}_0$}: $\varepsilon_0$ (same as PA)
\item \textbf{$\mathsf{ATR}_0$}: $\Gamma_0$ (Feferman-Schütte ordinal)
\item \textbf{$\Pi^1_1\text{-}\mathsf{CA}_0$}: $\psi_0(\Omega_\omega)$ (Takeuti-Feferman-Buchholz ordinal)
\item \textbf{ZFC}: Unknown; far beyond current ordinal analysis
\end{itemize}

\subsection{Natural Independence Results}
\label{sec:godel:natural}

The Gödel sentence $G$ is a logical artifact --- constructed specifically to be unprovable. For decades, logicians wondered whether \emph{natural} mathematical statements could be independent of PA. The answer came in the 1970s and 1980s.

\subsubsection*{Paris-Harrington Theorem (1977)}

Ramsey's theorem for graphs (finite version): For any $k, r, m$, there exists $N$ such that for any $r$-coloring of $k$-element subsets of $\{1, \dots, N\}$, there is a homogeneous set of size $m$ (a set all of whose $k$-subsets have the same color).

The Paris-Harrington strengthening adds a \emph{regulatory condition}: the homogeneous set must be \emph{large} in the sense that its size exceeds its smallest element. Formally:

\begin{theorem}[Paris-Harrington]
For any $k, r, m$, there exists $N$ such that for any $r$-coloring of $k$-element subsets of $\{1, \dots, N\}$, there is a homogeneous set $H$ of size at least $m$ with $|H| > \min(H)$.
\end{theorem}

This statement, denoted $\mathsf{PH}(k, r, m)$, is true in $\N$ but not provable in PA. The proof uses the infinite Ramsey theorem and a compactness argument (which requires induction beyond PA's reach). Paris and Harrington showed that $\mathsf{PH}$ implies the consistency of PA (when properly formalized), so by the second incompleteness theorem, $\mathsf{PH}$ is unprovable in PA.

More precisely: $\mathsf{PH}$ is equivalent (over PA) to $\mathsf{Con}(\mathsf{PA})$ plus some $\Pi_2$ reflection. This established the first ``natural'' independent statement for PA.

\subsubsection*{Goodstein's Theorem (1944, Kirby-Paris 1982)}

Write a natural number in \emph{hereditary base-$n$ notation}: express it in base $n$, then express the exponents in base $n$, and so on, until all numbers are below $n$. For example, 266 in base 2:

\[
266 = 2^8 + 2^3 + 2^1 = 2^{2^3} + 2^{2+1} + 2 = 2^{2^{2+1}} + 2^{2+1} + 2
\]

Now define the Goodstein sequence starting at $m$:
\begin{enumerate}
\item Write $m$ in hereditary base 2.
\item Replace every 2 with 3, then subtract 1.
\item Write the result in hereditary base 3.
\item Replace every 3 with 4, then subtract 1.
\item Continue indefinitely.
\end{enumerate}

\begin{theorem}[Goodstein's Theorem]
For any starting number $m$, the Goodstein sequence eventually reaches 0.
\end{theorem}

For example, starting with $m = 3$:
\begin{itemize}
\item Base 2: $3 = 2^1 + 1$
\item Replace 2 with 3: $3^1 + 1 = 4$, subtract 1: 3
\item Base 3: $3 = 3^1$
\item Replace 3 with 4: $4^1 = 4$, subtract 1: 3
\item Base 4: $3 = 3$ (no 4s), subtract 1: 2
\item Base 5: 2, subtract 1: 1
\item Base 6: 1, subtract 1: 0
\end{itemize}

The sequence reaches 0 after only a few steps. But for $m = 4$, the sequence grows astronomically before eventually descending. For $m = 19$, the sequence has length on the order of Graham's number $G_{64}$.

Goodstein's theorem is true in $\N$ but unprovable in PA. The proof uses transfinite induction up to $\varepsilon_0$: the ordinal associated with each term strictly decreases as the base increases, so the sequence must terminate. Since $\varepsilon_0$ induction is not available in PA, the theorem is not provable there.

Kirby and Paris (1982) proved that Goodstein's theorem implies the consistency of PA, establishing it as another natural independent statement.

%======================================================================
% SECTION 5: PHILOSOPHICAL CONSEQUENCES
%======================================================================
\section{Philosophical Consequences}
\label{sec:godel:consequences}

The incompleteness theorems are among the most consequential results in the intellectual history of mathematics. They did not merely solve a puzzle; they redrew the map of what mathematics, logic, and computation can and cannot do.

\subsection{What Became Possible}
\label{sec:godel:became_possible}

\begin{enumerate}
\item \textbf{Death of Hilbert's Program:} No finitary consistency proof for PA exists (since PA cannot prove its own consistency, and any finitary proof would be formalizable in PA). Hilbert's program survives only in a modified form: relative consistency proofs (e.g., Con(ZF) $\to$ Con(ZFC)) and proof-theoretic reductions.

\item \textbf{Independence Results:} Cohen (1963) used forcing to show $\mathsf{CH}$ (Continuum Hypothesis) is independent of ZFC. Gödel (1940) showed $\mathsf{CH}$ is consistent with ZFC. Incompleteness is not a bug --- it is a feature of rich theories. The independence phenomenon now encompasses hundreds of statements across set theory, combinatorics, and analysis.

\item \textbf{Undecidability of Logic:} Church (1936) and Turing (1936) used Gödel's methods to show first-order logic is undecidable (Entscheidungsproblem solved negatively). This followed directly from the encoding technique: if first-order logic were decidable, then PA's theorems would be decidable, contradicting incompleteness.

\item \textbf{Computability Theory:} The recursion theorem (Kleene 1938) generalizes the diagonal lemma: for any computable function $f$, there is an index $e$ such that $\varphi_e = \varphi_{f(e)}$. This is the basis of quines, viruses, and self-reproducing programs.

\item \textbf{Proof Theory:} Gentzen (1936) proved $\mathsf{Con}(\mathsf{PA})$ using transfinite induction up to $\varepsilon_0$. The \emph{proof-theoretic ordinal} measures a theory's strength. Modern proof theory continues to analyze stronger systems (from $\Pi^1_1$-CA$_0$ to subsystems of Z2).

\item \textbf{Algorithmic Information Theory:} Chaitin's $\Omega$ (1975) --- the halting probability --- is a real number that encodes the halting problem. Its bits are true but unprovable in any fixed formal system (incompleteness made quantitative). For any formal system $T$, only finitely many bits of $\Omega$ can be proven in $T$.

\item \textbf{Limits of Formal Verification:} Incompleteness imposes fundamental limits on automated theorem proving and program verification. No algorithm can decide all truths of arithmetic, and no formally verified system can prove its own correctness.
\end{enumerate}

\subsection{Natural Independence Results}
\label{sec:godel:natural_survey}

Gödel's original sentence $G$ is constructed specifically to be unprovable. But many ``natural'' mathematical statements are also independent:

\begin{itemize}
\item \textbf{Paris-Harrington Theorem (1977):} A strengthening of Ramsey's theorem (finite version) that is true but unprovable in PA.
\item \textbf{Goodstein's Theorem (1944, Kirby-Paris 1982):} Every Goodstein sequence eventually reaches 0. True but unprovable in PA.
\item \textbf{Kruskal's Tree Theorem (1960, Simpson 1985):} The set of finite trees is well-quasi-ordered under homeomorphic embedding. Unprovable in ATR$_0$.
\item \textbf{Robertson-Seymour Graph Minor Theorem (1985--2004):} Any minor-closed family of graphs is characterized by a finite set of forbidden minors. Proof uses well-quasi-ordering; unprovable in PA.
\item \textbf{Continuum Hypothesis (Gödel 1940, Cohen 1963):} Independent of ZFC.
\item \textbf{Whitehead's Problem (Shelah 1974):} Every $\aleph_1$-free abelian group is free. Independent of ZFC.
\item \textbf{Kaplansky's Conjectures (some independent):} Various conjectures in ring theory and group theory.
\item \textbf{Borel Determinacy (Martin 1975):} All Borel games are determined. Requires replacement (cannot be proved in Zermelo set theory).
\end{itemize}

\subsection{What the Theorems Taught Us}
\label{sec:godel:lessons}

\begin{enumerate}
\item \textbf{Truth Transcends Proof:} There are mathematical truths that no formal system can capture. Mathematical knowledge is not reducible to mechanical derivation. This is the \emph{Lucas-Penrose argument} (1961, 1989): human mathematicians ``see'' the truth of $G$, so human minds are not formal systems. (Controversial; see Feferman, Putnam.)

\item \textbf{The Inevitability of Self-Reference:} Any system rich enough to express arithmetic \emph{contains} self-reference. You cannot ban it (Russell tried). The diagonal lemma \emph{forces} it. Self-reference is not a paradox to avoid --- it is a structural feature of expressive systems.

\item \textbf{Consistency Cannot Be Self-Justified:} A system cannot prove its own consistency (unless it is inconsistent, in which case it proves everything). Trust in mathematics requires either a stronger system (which then needs \emph{its} consistency proved) or an act of faith (finitary intuition, or empirical success).

\item \textbf{The Map Is Not the Territory:} Formal systems are \emph{models} of mathematics, not mathematics itself. Incompleteness shows the model always leaves something out. This applies to \emph{any} formal model of a complex reality (physics, biology, law, AI alignment).

\item \textbf{Incompleteness as Generative:} Because no system is complete, mathematics is inexhaustible. New axioms (large cardinals, determinacy, forcing axioms) extend the frontier. Incompleteness is why mathematics never ends.
\end{enumerate}

\subsection{The Lucas-Penrose Argument: A Detailed Examination}
\label{sec:godel:lucas_penrose}

In 1961, philosopher John Lucas published ``Minds, Machines, and Gödel,'' arguing that Gödel's theorems prove the human mind is not a Turing machine. The argument was later revived and expanded by physicist Roger Penrose in \emph{The Emperor's New Mind} (1989) and \emph{Shadows of the Mind} (1994).

\textbf{The argument:}
\begin{enumerate}
\item Let $T$ be any formal system that captures the mathematical abilities of a human mind (or any Turing machine).
\item Gödel's first theorem gives a sentence $G_T$ that is true but unprovable in $T$.
\item A human mathematician can ``see'' that $G_T$ is true (by understanding the proof).
\item Therefore, human mathematicians can do something that no formal system $T$ can do --- they can recognize the truth of $G_T$.
\item Hence, the human mind is not a formal system (not a Turing machine).
\end{enumerate}

\textbf{Counterarguments and objections:}

\begin{itemize}
\item \textbf{Feferman's objection (systematic transcendence):} The human can only ``see'' the truth of $G_T$ by reasoning about \emph{the specific system $T$}. This reasoning uses a stronger system $T'$ (e.g., $T$ plus the assumption that $T$ is consistent). For any given formal system $T$, we can construct $G_T$ and a human can reason about it --- but this does not show the human transcends \emph{all} formal systems simultaneously. For any proposed formalization $M$ of human mathematical ability, $G_M$ is unprovable in $M$ but the human sees its truth by moving to $M'$ --- but this just shows the human's reasoning is not captured by $M$ in particular.

\item \textbf{Putnam's objection (consistency assumption):} To recognize $G_T$ as true, we must first recognize that $T$ is consistent. But recognizing $T$'s consistency may itself require transcending $T$. The human cannot ``see'' the truth of $G_T$ without already having reason to believe $T$ is consistent --- and that reason may be formalizable.

\item \textbf{Shapiro's objection (mechanism-compatible cognitive science):} Even if the mind is a formal system, the incompleteness theorems do not show it directly, because the mind's reasoning might be inconsistent. If the mind is inconsistent, it proves everything (including falsehoods), and Gödel's theorems do not apply.

\item \textbf{Chalmers' response:} The Lucas-Penrose argument fails because it cannot be formalized as a single effective procedure that outruns all formal systems. The human ability to ``see'' the truth of $G_T$ for each $T$ is itself a process that could be mechanized --- we just need to step through an enumeration of systems, but this would itself be a formal procedure.

\item \textbf{Boolos' observation:} There is no consistency requirement for the human mind. If the human mind is inconsistent, then it trivially proves $G$ (and also $\neg G$). Gödel's theorems only apply to consistent systems. An inconsistent mind would not be reliable, but Lucas-Penrose cannot claim the mind's consistency a priori.
\end{itemize}

Most philosophers and logicians today consider the Lucas-Penrose argument unsound. However, it continues to provoke valuable discussion about mechanism, consciousness, and the limits of computation.

\subsection{Mechanism vs. Platonism}
\label{sec:godel:mechanism_platonism}

The incompleteness theorems intersect with two major philosophical positions in the philosophy of mathematics:

\textbf{Mechanism:} The view that mathematical reasoning is algorithmic --- that the human mind, in its mathematical capacity, is equivalent to a Turing machine (or some formal system). Gödel's theorems challenge this view but do not refute it, because the human mind may be an inconsistent formal system, or may not be formalizable as a single system.

\textbf{Platonism:} The view that mathematical objects exist independently of human minds and that mathematical truth is discovered, not invented. Gödel was a committed Platonist, and he saw the incompleteness theorems as supporting Platonism: if truth transcends proof, then mathematical reality is richer than any formal system can capture. The independent existence of mathematical truth is exactly what incompleteness demonstrates.

Gödel wrote: ``I am convinced that the incompleteness theorems do not contradict Platonism, but rather support it. For if mathematical objects are real, we should expect them to have properties that we cannot fully capture by any finite system of axioms.''

\textbf{Formalism} (the view that mathematics is just symbol manipulation), by contrast, was severely weakened by incompleteness. If mathematics is just syntax, then the divergence of truth and provability is paradoxical. Platonism offers a natural explanation: there is a realm of mathematical objects with properties that no syntactic system can fully capture.

\textbf{Constructivism} and \textbf{intuitionism} were also affected. The incompleteness theorems apply to constructive systems as well (if they are rich enough to encode arithmetic). However, intuitionists reject the notion of ``truth in the standard model'' as a primitive, so the force of the theorems is different: they show the incompleteness of formal systems, not the transcendence of truth over proof.

Kurt Gödel was a mathematical Platonist. He believed that mathematical objects (sets, numbers, functions) exist independently of human thought and that the axioms of set theory describe a real universe. He wrote extensively on the philosophy of mathematics, arguing that intuition plays a role in discovering new axioms (such as large cardinal axioms) that extend the reach of formal systems. His Platonism influenced his work on the constructible universe $L$ and his belief that the Continuum Hypothesis has a definite truth value (unlike the view of formalists, who see it as a matter of choice).

\subsection{Modern Legacy: Where This Idea Appears Today}
\label{sec:godel:legacy}

\begin{itemize}
\item \textbf{Provability Logic (GL):} Modal logic of $\mathsf{Prov}$. $\Box \phi$ means ``$\phi$ is provable.'' Löb's Theorem: $\mathsf{PA} \vdash \mathsf{Prov}(\ulcorner \phi \urcorner) \to \phi$ iff $\mathsf{PA} \vdash \phi$. Fixed points in GL correspond to self-referential sentences. The logic GL is sound and complete for Kripke frames that are transitive and converse-well-founded.

\item \textbf{Independence in Set Theory:} Large cardinal axioms, $\mathsf{P} \neq \mathsf{NP}$ (possibly independent?), Whitehead's problem (Shelah 1974), Kaplansky's conjectures. The forcing technique (Cohen 1963) has been extended to produce independence results for almost every interesting statement in set theory.

\item \textbf{Automated Theorem Proving:} Incompleteness means no algorithm can decide all theorems. ATP systems (Lean, Coq, Isabelle) use human guidance + automation. The ``hammer'' tactics (Sledgehammer) call external provers. Gödel's theorems explain why ATP systems can never be fully automatic for sufficiently rich theories.

\item \textbf{Program Verification:} Verifying a program's correctness requires proving statements in a formal system. Incompleteness means some true properties are unprovable --- but we can often strengthen the system (add axioms, use richer logics). Practical verification works around incompleteness by restricting to decidable fragments (e.g., Presburger arithmetic, SMT solvers).

\item \textbf{AI and Formal Reasoning:} LLMs trained on formal libraries (Mathlib, Lean) can suggest proofs. But they cannot overcome incompleteness --- they just shift the boundary. The fundamental limitation remains: any fixed system has blind spots.

\item \textbf{Physics and Computation:} Wolfram's Principle of Computational Equivalence suggests physical systems exhibit Gödel-like incompleteness. The ``theory of everything'' may be formally incomplete. The Gödel metric (1949) provides an exact solution to Einstein's equations that allows closed timelike curves (time travel), reflecting Gödel's sustained engagement with physics.

\item \textbf{Cryptography and Zero-Knowledge Proofs:} Gödel's encoding of proofs as numbers anticipates the use of encodings in cryptography. Zero-knowledge proofs allow one to prove knowledge of a mathematical proof without revealing it --- a direct descendant of the arithmetization of metamathematics.
\end{itemize}

\subsection{Influence on Computer Science}
\label{sec:godel:cs}

The most concrete legacy of the incompleteness theorems is the transformation they enabled in computer science. Gödel numbering (encoding syntax as arithmetic) is the direct predecessor of Turing's encoding of machines as data. Both rely on the same principle: \emph{syntax can be made into data}. Turing read Gödel's 1931 paper and used the same technique for his 1936 paper. The stored-program computer (von Neumann architecture) is a physical realization of this principle: programs are data.

The direct lines of influence run through every chapter of this book:

\begin{itemize}
\item Chapter 1 (Computation): Gödel numbering $\to$ Turing's encoding of machines.
\item Chapter 2 (Halting): Incompleteness $\leftrightarrow$ Undecidability. $G$ is true $\leftrightarrow$ halting problem instance.
\item Chapter 4 (Rice): Semantic properties undecidable $\leftarrow$ incompleteness.
\item Chapter 7 (Kolmogorov): Chaitin's $\Omega$ --- incompleteness as randomness.
\item Chapter 11 (Zero Knowledge): Proving knowledge of a proof without revealing it.
\end{itemize}

Three further consequences deserve emphasis. First, the recursion theorem (Kleene 1938) generalizes the diagonal lemma: for any computable function $f$, there is an index $e$ such that $\varphi_e = \varphi_{f(e)}$. This is the basis of quines, viruses, and self-reproducing programs. Second, Chaitin's $\Omega$ (1975) --- the halting probability --- is a real number that encodes the halting problem; its bits are true but unprovable in any fixed formal system, a quantitative version of incompleteness. Third, incompleteness imposes a fundamental limit on formal verification: no algorithm can decide all truths of arithmetic, and no formally verified system can prove its own correctness. Practical verification works around this by restricting to decidable fragments, but the boundary Gödel drew is absolute.

%======================================================================
% SUMMARY
%======================================================================
\section{Summary}
\label{sec:godel:summary}

This chapter examined Gödel's incompleteness theorems, which demonstrate inherent limitations in all formal axiomatic systems capable of expressing elementary arithmetic. The first incompleteness theorem establishes that any consistent formal system $T$ extending Peano Arithmetic (PA) contains true statements that are unprovable within $T$. The proof constructs a Gödel sentence $G$ via the diagonal lemma, asserting its own unprovability: $\mathsf{PA} \vdash G \leftrightarrow \neg \mathsf{Prov}(\ulcorner G \urcorner)$. The second incompleteness theorem proves that no consistent extension of PA can prove its own consistency: $\mathsf{PA} \nvdash \mathsf{Con}(\mathsf{PA})$. The chapter detailed the arithmetization of syntax through Gödel numbering, the representability of primitive recursive predicates in PA, and the derivability conditions of Löb. Rosser's strengthening (1936) eliminated the assumption of $\omega$-consistency, relying only on simple consistency. The chapter also discussed Tarski's undefinability of truth, Löb's theorem, and natural independence results such as the Paris--Harrington theorem, Goodstein's theorem, and the Kirby--Paris result on Hydra games. The connection to computability theory was established: the set of theorems of PA is recursively enumerable but not recursive, and the undecidability of the halting problem provides an alternative proof of incompleteness. Several open problems were surveyed, including the search for natural independent statements, the proof-theoretic ordinal of PA ($\varepsilon_0$), and the consistency strength of stronger systems.

\section*{Key Takeaways}
\begin{itemize}
\item Gödel's First Incompleteness Theorem: any consistent formal system extending PA contains unprovable true statements, constructed via the diagonal lemma.
\item Gödel's Second Incompleteness Theorem: no consistent extension of PA can prove its own consistency.
\item The arithmetization of syntax (Gödel numbering) enables self-reference within arithmetic, transforming metamathematical claims into arithmetic statements.
\item Rosser's theorem strengthens the first incompleteness result by eliminating the need for $\omega$-consistency; Löb's theorem characterizes provability within formal systems.
\item Open problems include the classification of natural independent statements beyond PA, the computation of proof-theoretic ordinals for stronger theories, and the relationship between incompleteness and computational complexity.
\end{itemize}

%======================================================================
% EXERCISES
%======================================================================
\section{Exercises}
\label{sec:godel:exercises}

\begin{exercise}[Basic]
Prove that the set of Gödel numbers of theorems of PA is recursively enumerable but not recursive.
\end{exercise}

\begin{exercise}[Basic]
Show that if PA proves $\mathsf{Con}(\mathsf{PA})$, then PA is inconsistent.
\end{exercise}

\begin{exercise}[Basic]
Compute the Gödel number of the formula $S0 = S0$ using the encoding table from Section~\ref{sec:godel:numbering}.
\end{exercise}

\begin{exercise}[Basic]
Define a primitive recursive function for the factorial $n!$. Show that it is representable in PA.
\end{exercise}

\begin{exercise}[Intermediate]
Construct a Rosser sentence $R$ for PA: $R$ says ``For every proof of me, there is a shorter proof of $\neg R$.'' Prove that if PA is consistent, neither $R$ nor $\neg R$ is provable (without assuming $\omega$-consistency).
\end{exercise}

\begin{exercise}[Intermediate]
Prove that $\mathsf{Prov}$ is a $\Sigma_1$ formula: $\mathsf{Prov}(y) \equiv \exists x\, \mathsf{Proof}(x, y)$ where $\mathsf{Proof}$ is primitive recursive. Show that if $\phi$ is a $\Sigma_1$ sentence and true in $\N$, then $\mathsf{PA} \vdash \phi$ ($\Sigma_1$-completeness).
\end{exercise}

\begin{exercise}[Intermediate]
Derive the Second Incompleteness Theorem from Löb's Theorem (set $\phi = \bot$). Explain why this derivation relies on the derivability conditions.
\end{exercise}

\begin{exercise}[Intermediate]
Let $T$ be a consistent extension of PA. Show that there are infinitely many independent sentences in $T$ (sentences neither provable nor refutable).
\end{exercise}

\begin{exercise}[Intermediate]
Prove Tarski's Undefinability of Truth: there is no formula $\mathsf{True}(v)$ in the language of PA such that $\mathsf{PA} \vdash \mathsf{True}(\ulcorner \phi \urcorner) \leftrightarrow \phi$ for all sentences $\phi$.
\end{exercise}

\begin{exercise}[Intermediate]
Show that $\mathsf{PA} \vdash \phi$ implies $\mathsf{PA} \vdash \mathsf{Prov}(\ulcorner \phi \urcorner)$ (derivability condition 1). Why does this require that the $\mathsf{Proof}$ predicate is correctly defined?
\end{exercise}

\begin{exercise}[Advanced]
Prove Löb's Theorem in detail: assuming $\mathsf{PA} \vdash \mathsf{Prov}(\ulcorner \phi \urcorner) \to \phi$, show $\mathsf{PA} \vdash \phi$.
\end{exercise}

\begin{exercise}[Advanced]
Use the diagonal lemma to construct a sentence $G$ such that $\mathsf{PA} \vdash G \leftrightarrow \neg \mathsf{Prov}(\ulcorner G \urcorner)$. Verify the construction step by step.
\end{exercise}

\begin{exercise}[Advanced]
Show that the Paris-Harrington theorem is equivalent (over PA) to the consistency of PA plus some reflection. Explain why this implies Paris-Harrington is unprovable in PA.
\end{exercise}

\begin{exercise}[Challenge]
Research Gentzen's consistency proof for PA using transfinite induction up to $\varepsilon_0$. Explain why this does not violate the Second Incompleteness Theorem. What is the proof-theoretic ordinal of PA?
\end{exercise}

\begin{exercise}[Challenge]
Research \emph{ordinal analysis}: How are proof-theoretic ordinals computed for theories like $\mathsf{ACA}_0$, $\mathsf{ATR}_0$, $\Pi^1_1\text{-}\mathsf{CA}_0$? What is the ordinal of ZFC?
\end{exercise}

\begin{exercise}[Challenge]
Prove that Goodstein's theorem implies the consistency of PA using the ordinal assignment to Goodstein sequences. (Hint: assign $\varepsilon_0$ ordinals to each term and show they strictly decrease.)
\end{exercise}

\begin{exercise}[Computational]
Implement Gödel numbering for a simple formal system (e.g., propositional logic with axioms). Encode the statement ``This formula is not provable'' and verify it is undecidable in the system.
\end{exercise}

\begin{exercise}[Computational]
Write a program that enumerates theorems of a simple formal system and searches for a Gödel sentence. Verify that the search never terminates for the Gödel sentence.
\end{exercise}

\begin{exercise}[Computational]
Implement the Goodstein sequence function. Compute the length of the Goodstein sequence for $m = 3, 4, 5$. Observe the explosive growth for $m \ge 4$.
\end{exercise}

\begin{exercise}[Research]
Read Gödel's original 1931 paper (English translation). Identify the key differences between his presentation and the modern treatment using the diagonal lemma and primitive recursive functions. What role did $\omega$-consistency play in the original?
\end{exercise}

\begin{exercise}[Research]
Investigate the history of the Königsberg conference of 1930. Who attended? How did von Neumann react to Gödel's announcement? Write a short summary (250 words) of the event.
\end{exercise}

\begin{exercise}[Challenge]
Show that the set of $\Sigma_1$ truths (existential statements about $\N$) is recursively enumerable but not recursive. Deduce that there is no algorithm to decide $\Sigma_1$ truth, and explain the connection to the undecidability of the halting problem.
\end{exercise}

\begin{exercise}[Challenge]
Formulate the principle that ``provability implies truth'' as $\mathsf{Prov}(\ulcorner \phi \urcorner) \to \phi$. Show that PA cannot prove this principle for all $\phi$ (using Löb's theorem). This is the \emph{reflection principle} --- the system cannot prove its own soundness.
\end{exercise}

%======================================================================
% TIMELINE
%======================================================================
\section{Timeline}
\label{sec:godel:timeline}
\addcontentsline{toc}{section}{Timeline}

\begin{tabularx}{\linewidth}{|l|X|}
\hline
\textbf{Year} & \textbf{Event} \\ \hline
1900 & Hilbert's 23 problems; 2nd problem: consistency of arithmetic \\ \hline
1901 & Russell's Paradox \\ \hline
1906 & Gödel born in Brünn, Austria-Hungary \\ \hline
1908 & Russell's Type Theory; Zermelo's set theory \\ \hline
1910--13 & \emph{Principia Mathematica} (Whitehead \& Russell) \\ \hline
1924 & Gödel enters University of Vienna \\ \hline
1928 & Hilbert-Ackermann textbook poses completeness problem \\ \hline
1929 & Gödel's Completeness Theorem (first-order logic) \\ \hline
1930 & Gödel awarded PhD; Completeness published \\ \hline
1930 & Königsberg Conference; Gödel announces incompleteness \\ \hline
1931 & Gödel's Incompleteness Theorems published \\ \hline
1931 & Herbrand's theorem; general recursive functions (Gödel) \\ \hline
1933 & Gödel moves to Princeton; joins IAS \\ \hline
1936 & Church: $\lambda$-calculus undecidability; Turing: halting problem \\ \hline
1936 & Rosser: removes $\omega$-consistency requirement \\ \hline
1936 & Gentzen: consistency of PA via transfinite induction ($\varepsilon_0$) \\ \hline
1938 & Kleene: Recursion Theorem (generalized diagonalization) \\ \hline
1940 & Gödel: $\mathsf{CH}$ consistent with ZFC (constructible universe $L$) \\ \hline
1949 & Gödel metric: rotating universe solution to Einstein's equations \\ \hline
1955 & Feferman: Arithmetization of metamathematics \\ \hline
1963 & Cohen: Forcing; $\mathsf{CH}$ independent of ZFC \\ \hline
1974 & Paris-Harrington: First ``natural'' independent statement for PA \\ \hline
1975 & Chaitin: Algorithmic information theory; $\Omega$ number \\ \hline
1977 & Kirby-Paris: Goodstein's theorem independent of PA \\ \hline
1978 & Gödel dies in Princeton, New Jersey \\ \hline
1982 & Kirby-Paris: Goodstein's theorem independent of PA \\ \hline
1985 & Robertson-Seymour Graph Minor Theorem (unprovable in PA) \\ \hline
1989 & Penrose: \emph{The Emperor's New Mind} (Lucas-Penrose argument) \\ \hline
2001 & Feferman: analysis of Lucas-Penrose; definitive refutation \\ \hline
2000s & Proof assistants (Coq, Isabelle, Lean) formalize incompleteness \\ \hline
2010s & Homotopy Type Theory (Univalent Foundations) as new foundation \\ \hline
\end{tabularx}

%======================================================================
% CITATIONS
%======================================================================

%======================================================================
% INDEX ENTRIES
%======================================================================
\index{Gödel, Kurt}
\index{Incompleteness Theorems}
\index{Gödel numbering}
\index{diagonal lemma}
\index{self-reference}
\index{liar paradox}
\index{provability}
\index{consistency}
\index{formal system}
\index{Hilbert's program}
\index{Rosser's theorem}
\index{Gentzen}
\index{proof theory}
\index{independence results}
\index{forcing}
\index{Chaitin's Omega}
\index{proof-theoretic ordinal}
\index{transfinite induction}
\index{Löb's theorem}
\index{Paris-Harrington theorem}
\index{Goodstein's theorem}
\index{Tarski's undefinability of truth}
\index{primitive recursive functions}
\index{representability}
\index{derivability conditions}
\index{Hilbert-Bernays-Löb conditions}
\index{Lucas-Penrose argument}
\index{Platonism}
\index{Vienna Circle}
\index{Königsberg Conference}
\index{Einstein, Albert}
\index{ordinal analysis}
\index{recursion theorem}
\index{$\beta$-function}
\index{cut elimination}
